\documentclass[12pt]{article}
\usepackage{geometry}
\geometry{a4paper, portrait, margin=1in}

% --- Math Packages ---
\usepackage{amsmath}
\usepackage{amssymb}
\usepackage{amsfonts}
\usepackage{amsthm}

% --- Graphics & Diagrams ---
\usepackage{graphicx}
\usepackage{tikz}
\usetikzlibrary{shapes, arrows, positioning}

% --- Formatting ---
\usepackage{parskip}
\usepackage[hidelinks]{hyperref}
\usepackage{natbib}
\usepackage{enumitem}
\pagestyle{plain}

\title{MPhil Thesis Structure}
\author{Mohammad-Shah Karim}
\date{March 2026}

\begin{document}

\maketitle

\tableofcontents

\newpage

\section{Introduction}

\begin{itemize}
    \item Discuss pharmaceutical innovation
    \item Contributions in this paper - state the general problem / area of research.
    \item Part of the contribution is to collect and collate this data which is relevant for the question.
    \item Summary of Medicaid
    \item Discuss ACA: What changed in Medicaid, and what this means
    \item How is the ACA related to pharmaceutical innovation
    \item Why is it relevant?
    \item How will this paper be structured?
\end{itemize}

\section{Literature Review}

\begin{itemize}
    \item Literature on Medicaid / ACA - show that we're looking into literature that is relatively active. If they're new papers that works better. Want to give a feeling that when writing the paper we're speaking to an audience that is relatively large, and the audience looking at this policy is relatively large for different types of outcomes but this is just to show that the work is nested in an active research field.
    \item Public procurement and pharmaceutical pricing
    \item Market size and pharmaceutical innovation
    \item Dynamic innovation incentives
\end{itemize}

\section{The ACA}

\begin{itemize}
    \item PDL
    \item Rebates
    \item Market size
    \item Insurance
    \item Flag that we're not using staggered roll out to manage expectations
\end{itemize}


\section{Data}

\begin{itemize}
    \item SDUD
    \item Drugs@FDA: 
    \item Patents
    \item PDL
    \item RxNorm
\end{itemize}

\section{Methodology}

\begin{itemize}
    \item RQ1: Overall innovation response to market size. FEOLS, FE Poisson, Ex-ante, Ex-post
    \item RQ2: Disinentangling the effects of market sizing and pricing pressure. FEOLS, FE Poisson, Ex-ante, Ex-post
    \item RQ3: Innovation portfolio response - the DDD estimator.
    \item Calculation of Medicaid Share
    \item Calculation of PDL exposure
    \item Parallel trends tests
    \item Firm level variation
    \item PDL Counterfactual analysis - potential research question?
\end{itemize}

\section{Results}

\begin{itemize}
    \item Parallel trends
    \item RQ1
    \item RQ2
    \item RQ3 
\end{itemize}

\section{Theoretical Model}

\begin{itemize}
    \item Description of game
    \item Stage 1
    \item Stage 2
    \item Stage 3
    \item State dependent innovation - period 1 outcomes have implications for period 2 effort.
    \item Mapping theoretical states to empirical outcomes
    \item Equilibrium Analysis
    \item Comparative Statics
    \item Explaining the results
\end{itemize}

\section{Policy Implications}

How do you weigh the present benefits of consumer surplus with the cost to innovation in the future. As a social planner, planning for your rebate you need to account for this. Discussions should mostly look into this

\begin{itemize}
    \item Tension between states wanting lower net drug prices while firms need sufficiently high expected returns to justify R\&D spending
    \item Theoretical model highlights rebate mechanism is largely responsible for the current results
    \item Lower harm inflicted to firms through rebate mechanism
    \item Lower the cost of R\&D through subsidies or other incentives
\end{itemize}

\section{Extensions}

\begin{itemize}
    \item Currently map data to ATC, more granularity would be better.
    \item Clinical trial data
    \item PDL data from other states
    \item Theoretical model is illustrative, include an N player game with multiple winners
    \item Demand is assumed to be perfectly inelastic
    \item Logit model
\end{itemize}

\section{Conclusion}

Just a conclusion

\end{document}